\chapter{Background}
\label{ch:bg}

This chapter introduces the representation, renderer, prediction models, and
tracking system used by Splatt3R-SLAM. It also defines the metrics used in
later experiments.

\section{3D Gaussian splatting}
\label{sec:bg:3dgs}

\subsection{The representation}

3D Gaussian splatting~\cite{kerbl20233dgs} represents a scene as an unordered
set of anisotropic Gaussians. Primitive $n$ carries a centre
$\mu_n\in\mathbb{R}^3$, a covariance $\Sigma_n\in\mathbb{R}^{3\times3}$, an
opacity $\alpha_n\in(0,1)$ and a view-dependent colour encoded as
spherical-harmonic coefficients. Its unnormalised spatial influence at a
world point $x$ is
\begin{equation}
G_n(x) \;=\; \exp\!\Bigl(-\tfrac{1}{2}(x-\mu_n)^{\!\top}\Sigma_n^{-1}(x-\mu_n)\Bigr).
\label{eq:gaussian}
\end{equation}
A covariance must remain positive semi-definite under gradient descent.
It is therefore factored into a rotation
unit quaternion $q_n$ and positive per-axis scale $s_n$:
\begin{equation}
\Sigma_n \;=\; R(q_n)\,\operatorname{diag}(s_n^2)\,R(q_n)^\top .
\label{eq:cov}
\end{equation}
The scale is parameterised as $s_n=\exp(\hat{s}_n)$ elementwise, which keeps
it positive without a constraint. Large changes in $\hat{s}_n$ can produce the
scale excursions measured in \cref{sec:neg:lora}.

\subsection{Projection and rasterisation}

Rendering projects each Gaussian to the image plane. Let $R_{\mathcal CW}$ be
the rotation from world coordinates to the rendering camera and let $J_n$ be
the $2\times3$ Jacobian of perspective projection at the transformed centre.
The EWA construction of \citet{zwicker2001ewa} gives
\begin{equation}
\Sigma_n^{2D}
  =J_n R_{\mathcal CW}\Sigma_n^W R_{\mathcal CW}^{\top}J_n^{\top}.
\label{eq:proj}
\end{equation}
The rasteriser sorts primitives by depth per screen tile and composites them
front-to-back. Let $\mathbf c_\ell(u)$ be the SH-evaluated RGB colour and
$b_\ell(u)\in[0,1]$ the opacity multiplied by the 2D Gaussian footprint of the
$\ell$-th primitive along pixel ray $u$. The rendered colour is
\begin{equation}
\hat I(u)=\sum_{\ell=1}^{K}\mathbf c_\ell(u)b_\ell(u)
          \prod_{r<\ell}\bigl(1-b_r(u)\bigr),
\label{eq:alphacomp}
\end{equation}
and the residual transmittance after all $K$ primitives is
\begin{equation}
\mathcal T_{K}(u)=\prod_{\ell=1}^{K}\bigl(1-b_\ell(u)\bigr).
\label{eq:transmittance}
\end{equation}

\subsection{Background and visibility}
\label{sec:bg:black}

\Cref{eq:transmittance} is where the background enters. A renderer composites
$\hat I(u)\leftarrow\hat I(u)+\mathcal T_K(u)\mathbf c_{\mathrm{bg}}$, so with
$\mathbf c_{\mathrm{bg}}=0$ the residual
transmittance contributes \emph{zero radiance}: whatever a ray fails to hit is
black. For one isolated primitive with non-negative colour, increasing
opacity increases its contribution against that background. A stack is more
general: changing one opacity also changes the visibility of every primitive
behind it. A darker foreground primitive can hide a brighter background
primitive, so increasing its opacity can lower pixel luminance.
Opacity affects appearance as well as coverage, but the sign cannot be
inferred from the background alone. This distinction limits the
interpretation of the experiments in \cref{ch:thin}.

\subsection{Per-scene optimisation, and what a feed-forward model skips}

The reference 3DGS pipeline initialises primitives from a structure-from-motion
point cloud and optimises them against posed training images under
\begin{equation}
\mathcal{L}_{\text{3DGS}} = (1-\lambda_D)\,\lVert\hat{I}-I\rVert_1
      + \lambda_D\bigl(1-\mathrm{SSIM}(\hat{I},I)\bigr),
\label{eq:3dgsloss}
\end{equation}
interleaved with \emph{adaptive density control}: primitives whose view-space
positional gradient exceeds a threshold are cloned (if small) or split (if
large), primitives whose opacity falls below a threshold are pruned, and
opacity is periodically reset so that the optimiser must re-justify every
primitive it keeps. Spherical-harmonic degree is raised on a schedule, so that
low-frequency appearance is fitted before view-dependent effects.

A feed-forward model emits all primitives in one pass. It has no per-scene
densification, pruning schedule, or opacity reset. When predictions from many
keyframes are accumulated, the map therefore retains every inserted
primitive. \Cref{sec:lim:compact} measures the resulting size.
Our online refiner (\cref{ch:refiner}) restores the image objective in
\cref{eq:3dgsloss}, but not adaptive density control.

\section{Feed-forward pointmaps and Gaussians}
\label{sec:bg:feedforward}

\subsection{DUSt3R and MASt3R}

DUSt3R~\cite{wang2024dust3r} maps an uncalibrated image pair
$I_i,I_j\in\mathbb{R}^{H\times W\times3}$ to pointmaps
$X_{\mathcal C_i}^{(i)},X_{\mathcal C_i}^{(j)}
\in\mathbb{R}^{H\times W\times3}$ with confidences
$C_{\mathcal C_i}^{(i)},C_{\mathcal C_i}^{(j)}$. Here
$X_{\mathcal C_a}^{(b)}$ denotes image $b$'s pointmap expressed in camera
frame $\mathcal C_a$. Both images pass through a shared ViT encoder; two
decoders then attend to each other's tokens, so each view's prediction is
conditioned on the other, and DPT heads regress the pointmaps. The output lives
in the \emph{first} camera's frame, so a single forward pass already encodes
the relative pose implicitly rather than as an explicit estimate.

MASt3R~\cite{leroy2024mast3r} adds a matching head producing dense local
features $D^{(i)},D^{(j)}\in\mathbb{R}^{H\times W\times d}$, trained so that
corresponding pixels are mutual nearest neighbours in feature space. We write
$F_M(I_i,I_j)$ for one forward pass.

Two properties of this family govern everything downstream. The prediction is
\emph{scale-ambiguous up to a per-forward gauge} --- each independent pass
chooses its own scale, and nothing couples the choices --- and it is
\emph{pairwise}, so no single forward pass ever observes more than two views.
\Cref{sec:neg:frontend} measures scale disagreement and finds that joint
multi-view context reduces it in the tested models.

\subsection{Many-view feed-forward models}

VGGT~\cite{wang2025vggt} generalises the formulation: all views enter one
transformer with alternating frame-wise and global attention, and cameras,
depth maps and pointmaps are read out jointly. The consequence that matters
here is that a many-view model solves for a \emph{shared} scale across its
input window instead of drawing one per pair. \Cref{sec:neg:frontend}
quantifies the difference at $6.4\times$ --- and also shows that the same
weights run pairwise are worse than the two-view MASt3R line, so the advantage
belongs to the context, not to the architecture.

\subsection{Splatt3R}

Splatt3R~\cite{smart2024splatt3r} adds a third head, in parallel to the
pointmap and matching heads and using the same DPT decoder, which emits for
pixel $n$ a rotation quaternion $q_n\in\mathbb{R}^4$, a scale
$s_n\in\mathbb{R}_{>0}^3$, spherical-harmonic coefficients
$S_n\in\mathbb{R}^{3\times D_{\mathrm{SH}}}$, an opacity logit
$\widetilde\alpha_n\in\mathbb{R}$, and a mean offset
$\Delta_n\in\mathbb{R}^3$. The Gaussian centre and opacity are
\begin{equation}
\mu_n=x_n+\Delta_n,\quad x_n\in X_{\mathcal C_i}^{(i)},\qquad
\alpha_n=\operatorname{sigmoid}(\widetilde\alpha_n).
\label{eq:mu}
\end{equation}
$D_{\mathrm{SH}}$ is the number of SH basis coefficients per colour channel.
Quaternions are normalised and raw scale channels are exponentiated.
Offsets use a norm-dependent exponential rescaling of a signed vector, so their coordinates
are not restricted to positive values. In upstream training, \emph{only the Gaussian
head is trained}: the encoder and decoder remain those of the pre-trained
MASt3R model, and all Gaussian parameters are predicted in the first image's
camera frame, avoiding any ground-truth transformation at inference.

Covariance follows \cref{eq:cov} and rendering follows
\cref{eq:alphacomp}; the Gaussians a single forward pass emits are therefore
immediately renderable by a standard 3DGS rasteriser, with no optimisation
step in between.

Splatt3R is supervised only by novel-view renderings, under a mask $M$ of
pixels visible in at least one context view:
\begin{equation}
\mathcal{L}_{\mathrm{head}}
  =\mathcal{L}_{\mathrm{MSE},M}(\hat I,I)
  +0.25\,\mathcal{L}_{\mathrm{LPIPS(VGG)},M}(\hat I,I),
\label{eq:splatt3rloss}
\end{equation}
where the subscript $M$ denotes spatial averaging over valid pixels.
Two properties of \cref{eq:splatt3rloss} drive our method.
It is \emph{purely photometric} ---
depth enters only through the mask $M$ and through pair selection, never
through a gradient on geometry --- and it is \emph{pairwise}: the loss never
observes more than two context views, hence never observes a fused map.

In the released implementation the DC colour is anchored to the source image:
the head predicts a residual which is added to the spherical-harmonic encoding
of the input RGB value,
\begin{equation}
S_{\mathrm{DC}}(u)=S_{\mathrm{res}}(u)+\operatorname{RGB2SH}(I_i(u)).
\label{eq:colorresidual}
\end{equation}
Colour is clamped during conversion; it is not sigmoid-bounded. The residual
provides an appearance adjustment independently of opacity, so opacity is
not the only channel capable of compensating for input photometry.
Consequently, revisits can inherit exposure, white-balance, shadow, and
view-dependent differences from their source frames. Exposure normalisation
removes only the global component of this source-frame photometric gauge; it
does not enforce equality of $S_{\mathrm{DC}}$ for the same world surface.


\section{SLAM with a pointmap prior}
\label{sec:bg:slam}

We use $T_{AB}$ for a transform that maps coordinates from frame $B$ to
frame $A$. We use the front- and back-end of
MASt3R-SLAM~\cite{murai2025mast3rslam}. This section summarises the tracking
path whose trajectory accuracy is measured in \cref{ch:results}.

\subsection{Poses on the Sim(3) manifold}

Because network predictions are scale-ambiguous, poses are $\mathrm{Sim}(3)$
rather than $\mathrm{SE}(3)$, so that a per-keyframe scale is available to
absorb the ambiguity:
\begin{equation}
T=\begin{bmatrix} aR & t\\ 0 & 1\end{bmatrix},\qquad
T \leftarrow \tau \oplus T \triangleq \mathrm{Exp}(\tau)\circ T ,
\label{eq:sim3}
\end{equation}
with $a>0$ and $\tau\in\mathfrak{sim}(3)$. The scale degree of freedom
absorbs per-forward gauge error. This is why
\cref{sec:neg:frontend} finds that correcting the map's scales alone makes the
rendering worse.

\subsection{Matching and the ray error}

Under a generic central-camera assumption a pointmap induces its own camera
model through the ray map $\psi(X)=X/\lVert X\rVert$. Two rays differ by
\begin{equation}
\lVert\psi_1-\psi_2\rVert^2 = 2(1-\cos\theta), \qquad \cos\theta=\psi_1^{\!\top}\psi_2,
\label{eq:rayangle}
\end{equation}
which is \emph{bounded}, whereas a 3D point error is not. Since pointmap
predictions exhibit depth errors far more often than directional errors,
tracking minimises a ray error rather than a point error. For matches
$\mathcal M_{f,k}$ between the current frame $f$ and keyframe $k$,
\begin{equation}
E_r=\!\!\sum_{m,n\in\mathcal M_{f,k}}\!\!
\Bigl\lVert
\frac{\psi(\widetilde X_{\mathcal C_k,n}^{(k)})
 -\psi(T_{\mathcal C_k\mathcal C_f}X_{\mathcal C_f,m}^{(f)})}
{w(\omega_{m,n},\sigma_r^2)}
\Bigr\rVert_\rho ,
\label{eq:trackray}
\end{equation}
with $\lVert\cdot\rVert_\rho$ a Huber norm, $\omega_{m,n}$ the match
confidence, and $w$ its confidence-dependent noise scale; invalid matches
receive infinite scale. A small term on the difference
of distances from the camera centre keeps the pure-rotation case
non-degenerate. The system is solved by Gauss--Newton in an
iteratively-reweighted least-squares loop:
\begin{equation}
\bigl(J^{\!\top}WJ\bigr)\tau=-J^{\!\top}Wr,\qquad
T_{\mathcal C_k\mathcal C_f}\leftarrow
\tau\oplus T_{\mathcal C_k\mathcal C_f}.
\label{eq:gn}
\end{equation}
Both the iterative projection and the feature-space refinement of matches are
implemented as per-pixel CUDA kernels, which is what keeps tracking in the
low-millisecond range.

\subsection{Keyframing and canonical pointmap fusion}

A new keyframe is created when the number of valid matches, or of unique
keyframe pixels among them, falls below a threshold. Each keyframe maintains a
canonical pointmap updated by a confidence-weighted running average,
\begin{equation}
\begin{aligned}
\widetilde X_{\mathcal C_k}^{(k)}
&\leftarrow
\frac{\widetilde C_{\mathcal C_k}^{(k)}
      \widetilde X_{\mathcal C_k}^{(k)}
      +C_{\mathcal C_f}^{(k)}
       T_{\mathcal C_k\mathcal C_f}X_{\mathcal C_f}^{(k)}}
     {\widetilde C_{\mathcal C_k}^{(k)}+C_{\mathcal C_f}^{(k)}},\\
\widetilde C_{\mathcal C_k}^{(k)}
&\leftarrow\widetilde C_{\mathcal C_k}^{(k)}
             +C_{\mathcal C_f}^{(k)} .
\end{aligned}
\label{eq:fusion}
\end{equation}
Filtering merges information from many viewpoints into geometry that no single
small-baseline prediction could support. Note what \cref{eq:fusion} does
\emph{not} do: it averages within a keyframe, never across keyframes, so it
cannot remove a disagreement between two keyframes' scales. That disagreement
is the subject of \cref{sec:neg:frontend}.

\subsection{Loop closure and the back-end}

Loop-closure candidates are retrieved with an incremental adaptation of the
Aggregated Selective Match Kernel (ASMK)~\cite{tolias2016asmk} over encoded
features, as used by MASt3R-SfM. Each new keyframe queries the database for
top-$K$ candidates; those above a score threshold are passed to the decoder,
and bidirectional edges are added when the resulting match count clears a
second threshold. The keyframe's features then enter the inverted file index.
\Cref{sec:neg:retrieval} reports a closed attempt to refit these assets on our
own encoder's features.

\looseness=-1
Global consistency is obtained by minimising the same ray error over all edges
$\mathcal{E}$ of the keyframe graph,
\begin{equation}
E_g=\sum_{(i,j)\in\mathcal{E}}\ \sum_{m,n\in\mathcal M_{i,j}}
\Bigl\lVert
\frac{\psi(\widetilde X_{\mathcal C_i,m}^{(i)})
 -\psi(T_{\mathcal C_i\mathcal C_j}
       \widetilde X_{\mathcal C_j,n}^{(j)})}
{w(\omega_{m,n},\sigma_r^2)}
\Bigr\rVert_\rho,
\quad
T_{\mathcal C_i\mathcal C_j}=T_{W\mathcal C_i}^{-1}T_{W\mathcal C_j},
\label{eq:global}
\end{equation}
solved by second-order Gauss--Newton with sparse Cholesky over the $7N\times7N$
Hessian. The first $\mathrm{Sim}(3)$ pose is fixed to remove gauge freedom.
Splatt3R-SLAM does not change this solver; \cref{sec:exp:ate} compares its
trajectory with the inherited control.

\section{Metrics}
\label{sec:bg:metrics}

\subsection{Image metrics}

\looseness=-1
PSNR is a log-scaled mean-squared error,
$\mathrm{PSNR}=10\log_{10}(1/\mathrm{MSE})$ for images in $[0,1]$. It is
sensitive to global photometric offsets --- a uniform brightness shift costs
PSNR without changing structure --- which is why the protocol of
\cref{ch:protocol} normalises exposure identically for every system before
scoring. SSIM~\cite{wang2004ssim} compares local luminance, contrast and
structure statistics. LPIPS~\cite{zhang2018lpips} compares deep features of a
pre-trained network. Reconstruction evaluation uses LPIPS-AlexNet; head
adaptation retains the upstream VGG-based LPIPS loss.

\looseness=-1
We report PSNR and LPIPS together because they differ in
\cref{tab:replica}: PSNR improves on $8/8$ scenes and LPIPS on $5/8$.
LPIPS is not part of the refiner loss (\cref{eq:refloss}).

\subsection{Trajectory metrics}

\looseness=-1
Absolute Trajectory Error is translation RMSE after alignment to ground truth.
Monocular scale ambiguity requires a $\mathrm{Sim}(3)$ alignment, solved by
Umeyama's closed-form method~\cite{umeyama1991}. All five systems in
\cref{tab:ate} use \texttt{evo\_ape tum -as}.

\enlargethispage{\baselineskip}
\looseness=-1
Rendering also requires transporting ground-truth cameras into each map's
similarity gauge; \cref{ch:protocol} specifies this separate use of alignment.
